\section{System Requirements}

\subsection{Functional Requirements}

\subsubsection{User Account Management}

- The system must allow new users (readers) to register accounts by providing full name, email address, student or staff ID, and password. \\

- The system must authenticate users via username or email and password. \\

- The system must support at least three user roles: Reader, Librarian, and Administrator, each endowed with distinct authorization privileges.\\

- Users must be able to view and update their personal profile information, including contact details and notification preferences.\\

- The system must enable users to reset passwords via verification links dispatched to their registered email addresses.\\

- Administrators must be empowered to activate, deactivate, or temporarily suspend any user account.\\

- The system must log borrowing history for each reader, accessible to the respective reader and authorized librarians.\\



\subsubsection{Resource Catalog Management}

% Ghi chú: phần e-books, audiobooks, multimedia resources chưa được triển khai trong phiên bản này.

- The system must maintain a catalog of all library assets, including physical books and periodicals.\\

- Each catalog record must store: title, author(s), ISBN/ISSN, publisher, publication year, edition, genre/subject, language, physical location (shelf/classification code), total copy count, and available copy count.\\

- Librarians must be able to add, edit, and delete catalog records via the administration interface.\\

- The system must support batch catalog data ingestion via CSV and Excel (XLSX) file formats.\\

- The system must provide full-text search capability across title, author, ISBN, subject, and keyword fields.\\

- Search results must support filtering by category, language, availability status, publication year range, and media format.\\

- The system must display real-time copy availability for each item, including available, checked out, and reserved counts.\\



\subsubsection{Circulation Management}

- The system must allow readers to check out physical materials by barcode scanning or manual entry at circulation desks, or by placing online pickup requests.\\

- The system must enforce role-based loan limits (e.g., maximum 5 items for students, maximum 15 items for faculty).\\

- The system must assign due dates to borrowed items based on configurable loan duration rules per media type.\\

- The system must allow readers to renew checked-out items online up to a configurable maximum count, provided no hold requests exist for those items.\\

- The system must allow readers to place holds on items currently checked out by other users.\\

- The system must automatically notify the next user in a hold queue once a reserved item becomes available.\\

- The system must process item returns and immediately update copy availability in the catalog.\\

- The system must automatically calculate overdue fines based on daily configurable rates per material type.\\

- The system must enable readers to view active loans, due dates, holds, and outstanding fines from their user dashboard.\\

- The system must support inter-branch borrowing, enabling readers to request materials from other branch libraries in the network.\\



% \subsubsection{Digital Resource Access}

% - The system must provide secure access to e-books and digital journals.\\\

% - The system must enforce digital loan periods and automatically revoke access to digital files upon expiration.\\\

% - The system must support concurrent user access limits for digital assets (e.g., maximum N concurrent readers per licensed e-book).\\\





\subsubsection{AI Natural Language Search and Intelligent Interaction}

- The system must provide a natural language search interface allowing readers to query the catalog using conversational phrasing.\\

- The natural language search module must analyze user intent, extract attributes (genre, target audience, topic, length), and map them to catalog metadata.\\

- The system must display relevance-ranked natural language search results alongside standard keyword search outputs.\\

- The system must integrate a Virtual Assistant Chatbot for direct user interaction, resolving FAQ inquiries regarding library policies, operating hours, and loan procedures.\\



% \subsubsection{AI-Assisted Cataloging and Summarization}

% - Automatically auto-fill book metadata (Title, Author, Description) when librarians input an ISBN.\\



\subsubsection{Notifications and Communication}

- The system must dispatch automated automated email notifications for: 3-day pre-due-date reminders, overdue alerts, hold availability notices, and account status updates.\\

- Readers must be able to configure notification channels and frequencies within profile settings.\\

- The system must provide an in-app notification center displaying unread system alerts.\\

- Librarians must be able to broadcast mass announcements (e.g., library closure notices) to all users or target groups.\\



\subsubsection{Reporting and Analytics}

- The system must generate analytical reports regarding: top-borrowed items, overdue items, fine collection summaries, new user registrations, and catalog growth metrics.\\

- The system must feature a real-time dashboard for librarians displaying key operational metrics: active loans, overdue counts, copy availability, and daily transaction volumes.\\

- Reports must support export to PDF and Excel formats.\\



\subsubsection{System Administration and Configuration}

- Administrators must be able to configure global parameters including loan durations, fine rates, role borrowing limits, and hold expiration thresholds.\\

- The system must support multi-branch library management.\\

- The system must record full audit logs for administrative actions (configuration changes, account edits, batch ingestion).\\



\subsection{Non-Functional Requirements}



Non-functional requirements define the operational standards that ensure the system not only performs its business logic correctly, but also runs stably, securely, and delivers a positive user experience.



\subsubsection{Performance and User Experience}



\begin{itemize}

    % \item \textbf{Interface Responsiveness:} The user interface must respond smoothly to interactions. Common operations such as catalog browsing, book detail retrieval, and navigation between pages should feel immediate and fluid. For heavier operations --- such as AI-powered semantic search or bulk report generation --- the system provides visible loading indicators so users are informed of ongoing processing rather than perceiving the system as unresponsive.



    \item \textbf{Query Optimization:} Frequently accessed and infrequently changing data (book categories, branch locations, system configuration values) are cached at the application layer to reduce redundant database round-trips, thereby improving page rendering speed.



    \item \textbf{UI/UX Consistency:} Typography, color palette, navigation layouts, and interactive controls (Renew, Hold, Process Return) must remain visually and behaviorally consistent across all user roles and views, providing a coherent and professional user experience.

\end{itemize}



\subsubsection{Security and Access Control}



\begin{itemize}

    \item \textbf{Role-Based Access Control:} The system enforces strict RBAC combined with stateless JWT session tokens, ensuring that Readers, Librarians, and Administrators can only access interfaces and perform operations within their designated authorization boundaries.



    \item \textbf{Data Safety:} User passwords are never stored in plain text --- all credentials undergo salted hashing via bcrypt prior to persistence. All form inputs and API payloads are validated and sanitized server-side to prevent common injection and scripting vulnerabilities.



    % \item \textbf{Transport Security:} All client-server communications are conducted exclusively over TLS/HTTPS, protecting data in transit from interception or tampering.



    \item \textbf{Rate Limiting and CORS:} API endpoints are protected by rate-limiting middleware to mitigate brute-force attacks, and Cross-Origin Resource Sharing policies are explicitly configured to restrict unauthorized cross-domain access.

\end{itemize}



\subsubsection{Reliability and Maintainability}



\begin{itemize}

    \item \textbf{Data Consistency:} Critical business operations --- including loan record creation, fine assessment, and reservation queue updates --- are wrapped in ACID-compliant PostgreSQL transactions via Prisma ORM. This guarantees that no partial or inconsistent state is persisted in the event of mid-operation failures.



    \item \textbf{Code Quality and Maintainability:} The backend codebase is organized into clearly separated functional modules (authentication, catalog, circulation, notifications, AI service) with strict TypeScript static typing enforced throughout. This architecture limits the surface area for latent bugs and enables future developers to onboard, extend, and maintain individual modules independently without risk of cascading side-effects.



    \item \textbf{Recoverability:} The system supports on-demand database snapshots and point-in-time restoration of system configurations and branch data, providing a recovery mechanism against accidental data loss or configuration corruption.

\end{itemize}





